% !TeX root = ../informe.tex
% ===========================================================================
%  6. Reglas de negocio
%
%  Vale un 10 % de la rúbrica, «en especial la validación de solapamiento de
%  horarios». Lo que se evalúa es la trazabilidad: cada regla, dónde se
%  implementa y cómo se comprueba.
%
%  Los \label{rn:NN} de la tabla son los que hacen funcionar \rn{NN} desde
%  el resto del documento.
% ===========================================================================

\section{Reglas de negocio}
\label{sec:reglas-negocio}

\subsection{Catálogo de reglas}
\label{sec:catalogo-reglas}

\small
\begin{xltabular}{\textwidth}{C{1.6cm} Y Y}
  \caption{Reglas de negocio y dónde se implementa cada una.}
  \label{tab:reglas}\\
    \toprule
    \textbf{ID} & \textbf{Descripción} & \textbf{Implementación} \\
    \midrule
  \endfirsthead

    \multicolumn{3}{r}{\small\itshape Continuación de la Tabla \ref{tab:reglas}}\\
    \toprule
    \textbf{ID} & \textbf{Descripción} & \textbf{Implementación} \\
    \midrule
  \endhead

    \midrule
    \multicolumn{3}{r}{\small\itshape Continúa en la siguiente página}\\
  \endfoot

    \bottomrule
  \endlastfoot

    \textbf{RN-01}\label{rn:01} &
    Una reserva se realiza sobre una cancha, una fecha y un bloque horario
    predefinido de una hora. &
    \id{ms-reservas}: \id{BookingService.crear} crea el bloque mediante
    \id{TimeSlot.deUnaHoraDesde}. Prueba: \id{BookingServiceCreacionTest}. \\

    \textbf{RN-02}\label{rn:02} &
    No pueden existir dos reservas confirmadas que se solapen para la misma
    cancha. &
    \id{ms-reservas}: \id{BookingService.crear} realiza la comprobación previa
    y \id{BookingRepositoryAdapter} maneja el conflicto al persistir.
    \id{reservas_db}: restricción \id{reserva_sin_solape}.
    Pruebas: \id{BookingServiceCreacionTest} y \id{ConcurrenciaReservaIT}. \\

    \textbf{RN-03}\label{rn:03} &
    El usuario final solo puede cancelar sus propias reservas; el administrador
    puede cancelar cualquiera. &
    \id{ms-reservas}: \id{BookingService.cancelar} utiliza
    \id{X-User-Id} para comprobar el dueño de la reserva y
    \id{X-User-Role} para identificar al administrador.
    Prueba: \id{BookingServiceCancelacionTest}. \\

    \textbf{RN-04}\label{rn:04} &
    Una reserva no puede cancelarse cuando su fecha y hora de inicio ya
    ocurrieron, incluso si la operación la realiza un administrador. &
    \id{ms-reservas}: \id{BookingService.cancelar} utiliza
    \id{Booking.sePuedeCancelarEn}.
    Prueba: \id{BookingServiceCancelacionTest}. \\

    \textbf{RN-05}\label{rn:05} &
    Al cancelar una reserva, su estado cambia a \id{CANCELADA} y el bloque
    vuelve a quedar disponible sin eliminar el registro. &
    \id{ms-reservas}: \id{BookingService.cancelar} invoca
    \id{Booking.cancelar}. La restricción \id{reserva_sin_solape} solo
    considera reservas confirmadas.
    Prueba: \id{BookingServiceCancelacionTest}. \\

    \textbf{RN-06}\label{rn:06} &
    Un usuario final tiene un límite configurable de reservas activas
    simultáneas, con un valor de 3 por defecto. &
    \id{ms-reservas}: \id{BookingService.verificarTope} obtiene el valor
    mediante \id{ConfigurationRepositoryPort}.
    Prueba: \id{BookingServiceCreacionTest}. \\

    \textbf{RN-07}\label{rn:07} &
    Solo el administrador puede crear, editar o inactivar canchas y gestionar
    sus horarios y bloqueos de mantenimiento. &
    \id{ms-canchas}: \id{CourtService} aplica
    \id{exigirAdministrador} en las operaciones de gestión.
    Prueba: \id{CourtServiceTest}. \\

    \textbf{RN-08}\label{rn:08} &
    Una reserva puede encontrarse en estado \id{CONFIRMADA},
    \id{CANCELADA} o \id{FINALIZADA}. &
    \id{ms-reservas}: \id{Booking} nace como \id{CONFIRMADA},
    \id{Booking.cancelar} cambia el estado a \id{CANCELADA} y
    \id{Booking.estadoVigente} determina cuándo corresponde
    \id{FINALIZADA}. Pruebas: \id{BookingServiceCreacionTest} y
    \id{BookingServiceCancelacionTest}. \\

\end{xltabular}
\normalsize

\subsection{Validación de solapamiento de horarios}
\label{sec:solapamiento}

La \rn{02} establece que una misma cancha no puede tener dos reservas confirmadas
que se solapen en el mismo horario. El principal problema aparece cuando dos
usuarios intentan reservar el mismo bloque casi al mismo tiempo. En ese caso,
ambas solicitudes podrían comprobar la disponibilidad antes de que alguna de las
dos haya sido guardada y las dos podrían encontrar el horario libre.

Para evitar esta situación, la regla se controla en dos momentos. Primero,
\id{BookingService.crear} comprueba si el bloque solicitado se solapa con una
reserva confirmada existente. Para realizar esta comparación se utiliza
\id{TimeSlot.seSolapaCon}. Esta validación permite detectar el conflicto antes
de intentar guardar la reserva en los casos normales.

Sin embargo, esta comprobación no es suficiente cuando dos solicitudes se
procesan de forma concurrente. Las dos podrían realizar la validación al mismo
tiempo y obtener como resultado que el bloque está disponible. Por esta razón,
la comprobación final se mantiene también en PostgreSQL mediante la restricción
\id{reserva_sin_solape}.

\begin{lstlisting}[
  language=SQL,
  caption={Restricción que protege la \rn{02} en V1\_\_reservas.sql.},
  label={lst:exclude}
]
ALTER TABLE reserva ADD CONSTRAINT reserva_sin_solape
    EXCLUDE USING gist (
        cancha_id WITH =,
        tsrange(fecha + hora_inicio, fecha + hora_fin, '[)') WITH &&
    ) WHERE (estado = 'CONFIRMADA');
\end{lstlisting}

La restricción compara el identificador de la cancha y el rango horario de la
reserva. Si ya existe una reserva confirmada para la misma cancha cuyo horario
se solapa con el nuevo, PostgreSQL rechaza la operación. Así se evita que dos
solicitudes concurrentes terminen reservando el mismo bloque aunque ambas hayan
superado la comprobación previa de \id{BookingService}.

La definición también incluye dos condiciones importantes para el funcionamiento
de las reservas:

\begin{itemize}
    \item \id{WHERE (estado = 'CONFIRMADA')} hace que únicamente las reservas
    confirmadas ocupen un bloque. Cuando una reserva cambia a
    \id{CANCELADA}, deja de participar en esta restricción y el horario vuelve
    a quedar disponible sin eliminar el registro.

    \item El rango \id{'[)'} incluye la hora de inicio y excluye la hora de
    finalización. Esto permite, por ejemplo, registrar una reserva de
    10:00 a 11:00 y otra de 11:00 a 12:00 para la misma cancha, ya que son
    bloques consecutivos y no existe solapamiento entre ellos.
\end{itemize}

La restricción requiere la extensión \id{btree_gist}, ya que
\id{cancha_id} debe compararse por igualdad dentro del índice GiST. Este
requisito se prepara en \texttt{infra/postgres/init/02-extensiones.sql},
antes de aplicar la migración \id{V1__reservas.sql}.

Si el conflicto se detecta al momento de guardar, \id{BookingRepositoryAdapter}
traduce el error de persistencia a \id{SlotAlreadyBookedException}. La misma
excepción se utiliza cuando el solapamiento se detecta previamente en
\id{BookingService}. De esta forma, el conflicto se maneja de la misma manera
sin importar en cuál de las dos validaciones haya sido detectado, y la API
responde con HTTP 409.

Este caso también se comprueba mediante \id{ConcurrenciaReservaIT}, donde se
intentan registrar dos reservas sobre el mismo bloque de forma concurrente. La
prueba verifica que solamente una de las operaciones consiga registrar la
reserva y que la otra sea rechazada por solapamiento.

\subsection{Estados de una reserva}
\label{sec:estados-reserva}

La \rn{08} establece tres estados para una reserva: \id{CONFIRMADA},
\id{CANCELADA} y \id{FINALIZADA}. Toda reserva se crea inicialmente en estado
\id{CONFIRMADA}.

Una reserva confirmada puede cambiar a \id{CANCELADA} cuando se realiza una
cancelación antes de la hora de inicio. Esta operación se realiza mediante
\id{Booking.cancelar}, que modifica el estado y registra la fecha y hora en que
se produjo la cancelación. Una vez cancelada, la reserva permanece en ese estado.

El estado \id{FINALIZADA} se obtiene de una forma diferente. La reserva no se
actualiza en la base de datos cuando termina su horario. Al consultarla,
\id{Booking.estadoVigente} comprueba si continúa confirmada y si la hora de
finalización del bloque ya ocurrió. Cuando esto sucede, la reserva se presenta
como \id{FINALIZADA}.

Mientras una reserva ya comenzó pero todavía no termina, continúa apareciendo
como \id{CONFIRMADA}. Sin embargo, ya no puede cancelarse debido a la
\rn{04}.

Por lo tanto, las transiciones que maneja la aplicación son las siguientes:

\begin{itemize}
    \item \id{CONFIRMADA} $\rightarrow$ \id{CANCELADA}, cuando se cancela antes
    de la hora de inicio.
    \item \id{CONFIRMADA} $\rightarrow$ \id{FINALIZADA}, cuando el bloque ya
    terminó y el estado se calcula al consultar la reserva.
\end{itemize}

No existen transiciones desde \id{CANCELADA} ni desde \id{FINALIZADA} hacia
otro estado.

\subsection{Traducción de reglas a respuestas HTTP}
\label{sec:errores}

Las reglas de negocio se mantienen separadas de los códigos HTTP. Cuando una
operación no puede realizarse, la lógica de la aplicación genera una excepción
y el \id{ErrorHandler} del microservicio correspondiente se encarga de
convertirla en una respuesta HTTP. De esta manera, las clases del dominio no
necesitan conocer detalles propios de la capa web.

En \id{ms-reservas}, los conflictos relacionados con el estado actual de una
reserva se responden principalmente con \id{409 Conflict}. En cambio, cuando
una cancha existe pero no puede aceptar la reserva solicitada, se utiliza
\id{422 Unprocessable Entity}. Los permisos se manejan mediante
\id{403 Forbidden} y la ausencia de un recurso mediante \id{404 Not Found}.

\begin{table}[H]
  \centering
  \caption{Correspondencia entre violación de regla y respuesta HTTP.}
  \label{tab:errores}
  \begin{tabular}{C{1.6cm} C{2.4cm} L{8.6cm}}
    \toprule
    \textbf{Regla} & \textbf{HTTP} & \textbf{Excepción y significado} \\
    \midrule

    \rn{02} &
    \textbf{409} &
    \id{SlotAlreadyBookedException}. El bloque solicitado ya se encuentra
    ocupado por una reserva confirmada. Se utiliza tanto cuando el conflicto
    se detecta en \id{BookingService} como cuando se detecta al guardar por
    medio de la restricción \id{reserva_sin_solape}. \\

    \rn{03} &
    \textbf{403} &
    \id{ForbiddenOperationException}. Un usuario final intenta cancelar una
    reserva que pertenece a otro usuario. \\

    \rn{04} &
    \textbf{409} &
    \id{PastBookingCancellationException}. La fecha y hora de inicio de la
    reserva ya ocurrieron y, por lo tanto, ya no puede cancelarse. \\

    \rn{06} &
    \textbf{409} &
    \id{ActiveBookingLimitExceededException}. El usuario alcanzó el tope de 
    reservas activas configurado para el sistema. \\

    \rn{07} &
    \textbf{403} &
    \id{ForbiddenOperationException}. Un usuario que no es administrador
    intenta realizar una operación de gestión sobre las canchas, horarios o
    bloqueos de mantenimiento. \\

    \midrule

    -- &
    \textbf{404} &
    \id{NotFoundException}. La reserva, la cancha o el usuario solicitado no existe. \\

    -- &
    \textbf{422} &
    \id{CourtNotBookableException}. La cancha existe, pero no admite la reserva
    porque está inactiva, el bloque ya pasó, se encuentra fuera del horario de
    atención o coincide con un bloqueo de mantenimiento. \\

    -- &
    \textbf{400} &
    La petición contiene datos inválidos, no cumple las validaciones definidas
    para la entrada o contiene un valor que no corresponde con el tipo
    esperado. \\

    \bottomrule
  \end{tabular}
\end{table}

Las \rn{01}, \rn{05} y \rn{08} no generan una respuesta de error propia.
La \rn{01} se controla al construir el bloque horario y mediante las
validaciones de entrada. Por su parte, la \rn{05} y la \rn{08} describen el
comportamiento y los estados de una reserva, por lo que no representan por sí
mismas una operación inválida que deba traducirse a un código HTTP.
